\section{Introduction to Groups}

\subsection{Basic Axioms and Examples}

Let $G$ be a group.

\begin{problems}
    \item Determine which of the following binary operations are associative:
    \begin{problems}
        \item the operation $\star$ on $\z$ defined by $a \star b = a - b$
        \item the operation $\star$ on $\r$ defined by $a \star b = a + b + ab$
        \item the operation $\star$ on $\q$ defined by $a \star b = \dfrac{a + b}{5}$
        \item the operation $\star$ on $\z \times \z$ defined by 
        \[
        (a,b) \star (c,d) = (ad + bc, \, bd)
        \]
        \item the operation $\star$ on $\q - \{0\}$ defined by $a \star b = \dfrac{a}{b}$
    \end{problems}
    \begin{solalph}
        \item Not associative: $(2 - 0) - 2 \neq 2 - (0 -2)$.
        \item Associative:
        \begin{align*}
            (a \star b) \star c & = (a + b + ab) \star c \\
            & = (a + b + ab) + c + (a + b + ab)c \\
            & = a + b + ab + c + ac + bc + abc \\
            & = a + b + c + bc + ab + ac + abc \\
            & = a + (b + c + bc) + a(b + c + bc) \\
            & = a \star (b + c + bc) \\
            & = a \star (b \star c)
        \end{align*}
        \item Not associative: $(1 \star 0) \star 2 = 11/25, 1 \star (0 \star 2) = 7/25$.
        \item Associative:
        \begin{align*}
            [(a, b) \star (c, d)] \star (e, f) & = (ad + bc, bd) \star (e, f) \\
            & = ((ad + bc)f + bde, bdf) \\
            & = (adf + bcf + bde, bdf) \\
            & = (adf + b(cf + de), bdf) \\
            & = (a, b) \star (cf + de, df) \\
            & = (a, b) \star [(c, d) \star (e, f)]
        \end{align*}
        \item Not associative: $(1 \star 2) \star 3 = 1/6, 1 \star (2 \star 3) = 3/2$.
    \end{solalph}
    \item Decide which of the binary operations in the preceding exercise are commutative.
    \begin{solalph}
        \item Not commutative: $1 - 0 \neq 0 - 1$.
        \item Commutative: $a \star b = a + b + ab = b + a + ba = b \star a$.
        \item Commutative: $a \star b = (a + b)/5 = (b + a)/5 = b \star a$.
        \item Commutative: $(a, b) \star (c, d) = (ad + bc, bd) = (cb + da, db) = (c, d) \star (a, b)$.
        \item Not commutative: $1/2 \neq 2/1$.
    \end{solalph}
    \item Prove that addition of residue classes in $\intmod$ is associative (you may assume it is well defined).
    \begin{sol}
        Let $\widebar a, \widebar b, \widebar c \in \intmod$. Then $(\widebar a + \widebar b) + \widebar c = \widebar{a + b} + \widebar c = \widebar{(a + b) + c} = \widebar{a + (b + c)} = \widebar a + \widebar{b + c} = \widebar a + (\widebar b + \widebar c)$.
    \end{sol}
    \item Prove that multiplication of residue classes in $\intmod$ is associative (you may assume it is well defined).
    \begin{sol}
        Uses similar reasoning as the previous exercise (where $\cdot$ is associative over $\z$).
    \end{sol}
    \item Prove for all $n > 1$ that $\intmod$ is not a group under multiplication of residue classes.
    \begin{sol}
        Note that in any $\intmod$, $\widebar 1$ is a multiplicative identity. Moreover, there is no $x \in \z$ such that $0x = 1$ so that $\widebar 0 \in \intmod$ has no multiplicative inverse.  
    \end{sol}
    \item Determine which of the following sets are groups under addition:
    \begin{problems}
        \item the set of rational numbers (including $0 = 0/1$) in lowest terms whose denominators are odd
        \item the set of rational numbers (including $0 = 0/1$) in lowest terms whose denominators are even
        \item the set of rational numbers of absolute value $< 1$
        \item the set of rational numbers of absolute value $\ge 1$ together with $0$
        \item the set of rational numbers with denominators equal to $1$ or $2$
        \item the set of rational numbers with denominators equal to $1$, $2$, or $3$
    \end{problems}
    \begin{sol}
        Let $S$ denote the set in each part.
        \begin{problems}
            \item Clearly, $0 \in S$ and for any $r \in S$, then $-r \in S$ is the additive inverse. To prove closure, let $a/b, c/d \in S$, and consider its sum $(ad + bc)/bd$. Since $b, d$ are odd, then $bd$ is also odd. Since having an even factor results in an even number, then any divisor of $bd$ must also be odd. Then $(ad + bc, bd)$ must also be odd, and if we divide both the numerator and the denominator by this quantity, then we still result in odd numbers. Hence, $A$ is closed under addition and is a group. 
            \item $S$ is not a group, since $1/2 \in S$, but $1/2 + 1/2 = 2/2 = 1/1 \not\in S$.
            \item Same as previous.
            \item $S$ is not a group, since $3/2, 1 \in S$, but $3/2-1 = 1/2 \not\in S$.
            \item Using the Division Algorithm, any $p \in S$ may be split into the form $a + b/2$, where $a$ is even and $b$ is either $0$ (if the numerator $p$ is even) or $1$ (if the numerator of $p$ is odd). Then if we have $a_1 + b_1/2, a_2 + b_2/2 \in S$ and we consider their sum, then either both $b_i$ is $0$ or $1$, in which case we remain an integer with denominator $1$, or just one of the $b_i$ is $1$ in which case we end with a rational with denominator $2$ (WLOG, suppose $b_1 = 1$ and $b_2 = 0$. Then we get the rational $(2a_1 + 2a_2 + 1)/2$). Hence, $S$ is closed under addition, and it also has identity $0$ and additive inverse $-p$ so that $S$ is a group.
            \item $S$ is not a group, since $1/2, 1/3 \in S$, but $1/2 + 1/3 = 5/6 \not\in S$. \qh
        \end{problems}
    \end{sol}
    \item Let $G = \{ x \in \r \mid 0 \le x < 1 \}$ and for $x, y \in G$ let $x \star y$ be the fractional part of $x + y$ (i.e., $x \star y = x + y - [x + y]$ where $[a]$ is the greatest integer less than or equal to $a$). Prove that $\star$ is a well defined binary operation on $G$ and that $G$ is an abelian group under $\star$ (called the \textit{real numbers mod $I$}).
    \begin{sol}
        Let $x, y \in G$. Then we have two cases:
        \begin{itemize}
            \item $0 \leq x + y < 1$: Then $[x + y] = 0$, and $x \star y = x + y \in G$,
            \item $1 \leq x + y < 2$: Then $[x + y] = 1$, and $x \star y = x + y - 1 \in G$.
        \end{itemize}
        so that $\star$ is well defined on $G$. To show associativity, let $z \in G$. Then there are three cases:
        \begin{itemize}
            \item If $0 \leq x + y, y + z < 1$, then $[x + y] = [y + z] = 0$:
            \begin{align*}
                (x \star y) \star z & = (x + y - [x + y]) \star z \\
                & = x + y + z - [x + y + z] \\
                & = x \star (y + z - [y + z]) \\
                & = x \star (y \star z)
            \end{align*}
            \item If $1 \leq  x + y, y + z < 2$, then $[x + y] = [y + z] = 1$:
            \begin{align*}
                (x \star y) \star z & = (x + y - [x + y]) \star z \\
                & = (x + y - 1) + z - [x + y - 1 + z] \\
                & = x + (y + z - 1) - [x + y + z - 1] \\
                & = x \star (y + z - 1) - [x + (y + z - 1)] \\
                & = x \star (y + z - [y + z]) \\
                & = x \star (y \star z)
            \end{align*}
            \item If $1 \leq x + y  < 2$ and $0 \leq y + z < 1$, then $[x + y] = 1$ and $[y + z] = 0$. The case of $0 \leq x + y < 1$ and $1 \leq y + z < 2$ is similar:
            \begin{align*}
                (x \star y) \star z & = (x + y - [x + y]) \star z \\
                & = (x + y - 1) + z - [x + y - 1 + z] \\
                & = x + y - 1 + z - ([x + y + z] - 1) \\
                & = x + y + z - 1 - [x + y + z] + 1 \\
                & = x + y + z - [x + y + z] \\
                & = x \star (y + z - [y + z]) \\
                & = x \star (y \star z)
            \end{align*}
        \end{itemize}
        Then $\star$ is associative. Moreover, commutativity of regular addition implies commutativity of $G$. For the identity element, note that $0 \in G$ and $0 \star x = x + 0 - [x] = x$. To determine the inverse of some $x \in G$, suppose $w \in G$ such that $x \star w = 0$. Then $x + w - [x + w] = 0$. Note that $0 \leq x + w < 2$, so $[x + w] = 0$ or $1$. If $[x + w] = 0$, then $w = -x$. Since $w, x \in G$, then this is when $x = w = 0$. If $[x + w] = 1$, then $w = 1 - x$ for $0 < x < 1$. Hence, $G$ is closed under inverses and is a group.
    \end{sol}
    \item Let $G = \{z \in \c :  z^n = 1 \text{ for some } n \in \zp \}$.
    \begin{problems}
        \item Prove that $G$ is a group under multiplication (called the group of roots of unity in $\c$).
        \item Prove that $G$ is not a group under addition.
    \end{problems}
    \begin{solalph}
        \item Note that $1 \in G$, so $G$ has an identity. Moreover, if $zz\inv = 1$ for some $z \in G$, then $z \inv = 1/z \in G$ because $(1/z)^n = 1/z^n = 1$. Lastly, let $z, w \in G$, where $z^n = 1$ and $w^m = 1$ for $n, m \in \zp$. Then
        $$(zw)^{mn} = (z^n)^m (w^m)^n = 1$$
        so that $zw \in G$. Moreover, multiplication is associative in $G$ since $G \subseteq \c$.
        \item Note that $1, -1 \in G$, but $1 + (-1) = 0 \not\in G$, since there is no $n$ such that $0^n = 1$.
    \end{solalph}
\end{problems}